\documentclass{beamer}
\usetheme{default}
\usecolortheme{seahorse}

\usepackage{amsmath,amssymb,mathrsfs,bm}
\usepackage{anyfontsize}
\usepackage{unicode-math}
\unimathsetup{math-style=ISO, bold-style=ISO, mathrm=sym}
\setmathfont{XITS Math}
\usepackage{fontspec}
\setsansfont{Libertinus Sans}

\AtBeginDocument{
  \let\uglyepsilon\epsilon
  \renewcommand\epsilon{\varepsilon}
}
\AtBeginDocument{
  \let\uglymathsf\mathsf
  \renewcommand\mathsf{\symsf}
}
\renewcommand{\bm}{\symbf}

\usepackage{color}
\usepackage{graphicx,caption,subcaption}
\usepackage{hyperref}
\hypersetup{
    colorlinks=true,
    linkcolor=blue,
    filecolor=green,
    urlcolor=cyan,
    citecolor=cyan,
}

\title{\large A natural boundary of dark matter haloes revealed around the minimum bias and maximum infall locations}
\author{Matthew Fong, Jiaxin Han \texorpdfstring{\\ \vspace{2em}}{} Speaker: Zhijun Zhang}
\date{2021.12.23}

\begin{document}

\begin{frame}[plain]
    \maketitle
\end{frame}

\begin{frame}{Introduction}
    \begin{itemize}
        \item[\bullet] Dark matter haloes are the basic structure in the Universe
        \item[\bullet] The evolution of dark matter haloes determines galaxy evolutions and properties
        \item[\bullet] The understanding of what is the size of a halo is essential for studying dark matter halo evolution
    \end{itemize}
\end{frame}

\begin{frame}{Dark Matter Halo Radius}
    \begin{itemize}
        \item[\bullet] Virial Radius $r_{\mathrm{vir}}$ (spherical and static)
        \item[\bullet] Splashback Radius $r_{\mathrm{sp}}$ (density profile reaches its deepest slope)
        \item[\bullet] Turnaround Radius $r_{\mathrm{ta}}$ (Hubble flow velocity overcomes peculiar radial velocity)
        \item[\bullet] $r_{\mathrm{vir}}\lesssim r_{\mathrm{sp}}<r_{\mathrm{ta}}$
    \end{itemize}
\end{frame}

\begin{frame}{Characteristic Depletion Radius}
    \begin{itemize}
        \item[\bullet] CosmicGrowth Simulations (high resolution N body simulation with $3072^3$ dark-matter particles)
        \item[\bullet] Bias profile (overdensity profile around a halo particle relative to average)
        \item[\bullet] Characteristic Depletion Radius $r_{\mathrm{cd}}$ (trough of the bias profile)
        \item[\bullet] Correlation between matter and haloes is the weakest
        \item[\bullet] Depends strongly on mass, environment and time, deeper trough = older
    \end{itemize}
\end{frame}

\begin{frame}{Characteristic Depletion Radius}
    \begin{figure}
        \includegraphics[width=1\linewidth]{Figure_1.jpg}
    \end{figure}
\end{frame}

\begin{frame}{Characteristic Depletion Radius}
    \begin{itemize}
        \item[\bullet] 2-3 times the splashback radius
        \item[\bullet] 2.5 times the virial radius (mass independent)
        \item[\bullet] approximately the turnaround radius
        \item[\bullet] average density in $r_{\mathrm{cd}}$ is about 40 times the background
    \end{itemize}
\end{frame}

\begin{frame}{Characteristic Depletion Radius}
    \begin{figure}
        \includegraphics[width=0.75\linewidth]{Figure_2.jpg}
    \end{figure}
\end{frame}

\begin{frame}{Inner Depletion Radius}
    \begin{itemize}
        \item[\bullet] Inner Depletion Radius $r_{\mathrm{id}}$ (maximum mass inflow rate)
        \item[\bullet] Smaller than $r_{\mathrm{cd}}$ by $10\sim20\%$
        \item[\bullet] Matter drained from outside $r_{\mathrm{id}}$ leading the trough at $r_{\mathrm{cd}}$
    \end{itemize}
\end{frame}

\begin{frame}{Characteristic Depletion Radius}
    \begin{itemize}
        \item[\bullet] Natural boundary of splashback particles
    \end{itemize}
    \begin{figure}
        \includegraphics[width=0.5\linewidth]{Figure_3.jpg}
    \end{figure}
\end{frame}

\begin{frame}{Characteristic Depletion Radius}
    \begin{itemize}
        \item[\bullet] Perfect radius of exclusion radius
    \end{itemize}
    \begin{figure}
        \includegraphics[width=0.5\linewidth]{Figure_4.jpg}
    \end{figure}
\end{frame}

\begin{frame}
    \begin{center}
        \Huge Thank you for listening!
    \end{center}
\end{frame}

\end{document}
